\documentclass[11pt, a4paper]{article}

% ── Packages ──────────────────────────────────────────────────────────────
\usepackage[margin=1in]{geometry}
\usepackage[utf8]{inputenc}
\usepackage[T1]{fontenc}
\usepackage{lmodern}
\usepackage{microtype}
\usepackage{graphicx}
\usepackage{booktabs}
\usepackage{array}
\usepackage{tabularx}
\usepackage{multirow}
\usepackage{enumitem}
\usepackage{xcolor}
\usepackage{hyperref}
\usepackage{titlesec}
\usepackage{fancyhdr}
\usepackage{amsmath}
\usepackage{caption}
\usepackage{float}

% ── Colors ────────────────────────────────────────────────────────────────
\definecolor{secblue}{RGB}{0, 51, 102}
\definecolor{linkblue}{RGB}{0, 76, 153}
\definecolor{lightgray}{RGB}{245, 245, 245}

% ── Hyperref setup ────────────────────────────────────────────────────────
\hypersetup{
    colorlinks=true,
    linkcolor=secblue,
    urlcolor=linkblue,
    citecolor=secblue
}

% ── Section formatting ────────────────────────────────────────────────────
\titleformat{\section}{\large\bfseries\color{secblue}}{\thesection.}{0.5em}{}
\titleformat{\subsection}{\normalsize\bfseries\color{secblue}}{\thesubsection}{0.5em}{}
\titlespacing*{\section}{0pt}{14pt}{6pt}
\titlespacing*{\subsection}{0pt}{10pt}{4pt}

% ── Header / Footer ──────────────────────────────────────────────────────
\pagestyle{fancy}
\fancyhf{}
\fancyhead[L]{\small\color{gray}Independent Study Report}
\fancyhead[R]{\small\color{gray}Spring 2026}
\fancyfoot[C]{\small\thepage}
\renewcommand{\headrulewidth}{0.4pt}

% ── Custom column type ───────────────────────────────────────────────────
\newcolumntype{L}[1]{>{\raggedright\arraybackslash}p{#1}}
\newcolumntype{C}[1]{>{\centering\arraybackslash}p{#1}}
\newcolumntype{R}[1]{>{\raggedleft\arraybackslash}p{#1}}

% ══════════════════════════════════════════════════════════════════════════
\begin{document}

% ── Title Block ───────────────────────────────────────────────────────────
\begin{center}
    {\LARGE\bfseries\color{secblue} Agentic Threat Detection \& Triage\\[4pt]
    from Network Telemetry}\\[16pt]
    {\large Uday}\\[4pt]
    {\normalsize Independent Study\,---\,Spring 2026}\\[4pt]
\end{center}

\vspace{-4pt}
\noindent\rule{\textwidth}{0.8pt}
\vspace{6pt}

% ══════════════════════════════════════════════════════════════════════════
\section{Aim}

To design and implement a \textbf{multi-agent, rule-based pipeline} that automatically
ingests raw network telemetry (Zeek and Suricata logs), detects threats through
multi-signal behavioral analysis, assembles structured investigation cases with
scored evidence, and validates each case through a critic feedback loop---reducing
the manual triage burden faced by Security Operations Center (SOC) analysts.

% ══════════════════════════════════════════════════════════════════════════
\section{Objectives}

\begin{enumerate}[leftmargin=*, itemsep=2pt]
    \item Build multi-signal detectors for \textbf{reconnaissance scanning} (3 signals) and \textbf{DNS beaconing} (4 signals) that go beyond signature matching.
    \item Implement a \textbf{4-agent orchestration pipeline} (Triage, Evidence, Critic, Report) that compresses raw detections into analyst-ready cases.
    \item Design a transparent \textbf{5-factor confidence model} so every score is explainable.
    \item Evaluate rigorously with a \textbf{14-PCAP benchmark} (12 synthetic + 2 public), a \textbf{5-configuration ablation study}, and \textbf{3 adversarial evasion tests}.
    \item Produce \textbf{17 diagnostic visualizations} and quantify results with Hedge's~$g$ effect size.
\end{enumerate}

% ══════════════════════════════════════════════════════════════════════════
\section{Data Used}

\subsection{Synthetic PCAPs (12 files, generated with Scapy)}

All synthetic PCAPs are deterministic and reproducible via a single script.

\begin{table}[H]
\centering\small
\caption{Synthetic PCAP suite}
\label{tab:pcaps}
\begin{tabularx}{\textwidth}{@{} l c L{6.5cm} @{}}
\toprule
\textbf{PCAP Name} & \textbf{Label} & \textbf{What It Simulates} \\
\midrule
synthetic-scan          & Malicious & Single source scanning 130+ destinations \\
synthetic-beacon        & Malicious & DNS C2 beaconing every 30\,s \\
synthetic-benign        & Benign    & Normal web browsing and DNS lookups \\
synthetic-mixed         & Malicious & Normal traffic with embedded port scan \\
synthetic-dns-exfil     & Malicious & High-entropy subdomain DNS exfiltration \\
synthetic-multi-attacker& Malicious & 3 simultaneous threats (scan + RDP + DNS) \\
synthetic-killchain     & Malicious & Recon $\rightarrow$ C2 $\rightarrow$ lateral movement \\
synthetic-slow-scan     & Malicious & Low-and-slow scan (1 conn / 30--60\,s) \\
synthetic-noisy-benign  & Benign    & Heavy benign traffic (false-positive stress test) \\
evasion-ip-rotation     & Malicious & 5 IPs each below fan-out threshold \\
evasion-jittery-beacon  & Malicious & DNS beacon with high timing jitter (CV\,>\,1.0) \\
evasion-slow-drip       & Malicious & 1 DNS query per 5\,min over 1 hour \\
\bottomrule
\end{tabularx}
\end{table}

\subsection{Public Datasets (2 files)}
\begin{itemize}[itemsep=2pt]
    \item \textbf{CTU-13 Neris Botnet} (56\,MB)\,---\,Real botnet traffic from the Czech Technical University dataset. \textit{Note:} Processing timed out on this capture; included for completeness and documented as a known limitation.
    \item \textbf{Wireshark HTTP Sample}\,---\,Benign HTTP traffic used as a clean baseline.
\end{itemize}

% ══════════════════════════════════════════════════════════════════════════
\section{Data Processing \& Normalization}

Raw PCAPs are processed through \textbf{Zeek} and \textbf{Suricata} running in Docker containers. Their heterogeneous outputs are parsed and merged into a unified \textbf{13-field normalized schema}:

\begin{center}\small
\texttt{ts \,|\, src\_ip \,|\, dst\_ip \,|\, src\_port \,|\, dst\_port \,|\, proto \,|\, event\_type \,|\, sensor \,|\, signature \,|\, severity \,|\, conn\_state \,|\, duration \,|\, metadata}
\end{center}

\noindent Key processing steps:
\begin{enumerate}[itemsep=2pt]
    \item \textbf{Zeek parsing}\,---\,Extract connection records (\texttt{conn.log}) and DNS queries (\texttt{dns.log}).
    \item \textbf{Suricata parsing}\,---\,Extract alert events from \texttt{eve.json}.
    \item \textbf{Schema unification}\,---\,Map both sensor outputs to the 13-field schema, preserving sensor-specific metadata as a JSON field.
    \item \textbf{Storage}\,---\,Normalized events stored as Parquet for efficient downstream analysis.
\end{enumerate}

% ══════════════════════════════════════════════════════════════════════════
\section{System Architecture}

The pipeline has two main stages: \textbf{baseline detection} and \textbf{agent orchestration}.

\subsection{Stage 1 --- Multi-Signal Baseline Detection}

Two detectors operate on normalized events using configurable thresholds:

\begin{table}[H]
\centering\small
\caption{Detector configurations}
\label{tab:detectors}
\begin{tabular}{@{} l l c c @{}}
\toprule
\textbf{Detector} & \textbf{Signal} & \textbf{Weight} & \textbf{Threshold} \\
\midrule
\multirow{3}{*}{Recon Scanning}
  & Fan-out (unique dest.\ IPs)      & 0.50 & 15 IPs \\
  & Burst (conns / second)            & 0.25 & 20 c/s \\
  & Failed connection ratio            & 0.25 & 0.40 \\
\midrule
\multirow{4}{*}{DNS Beaconing}
  & Repeated queries (same domain)     & 0.35 & 5 queries \\
  & Periodicity (CV of intervals)      & 0.30 & --- \\
  & NXDOMAIN ratio                     & 0.20 & 0.15 \\
  & Domain diversity                   & 0.15 & 3 domains \\
\bottomrule
\end{tabular}
\end{table}

\subsection{Stage 2 --- 4-Agent Orchestration Pipeline}

\begin{enumerate}[itemsep=4pt]
    \item \textbf{Triage Agent}\,---\,Groups detections by \texttt{(type, src\_ip, time\_bucket)} into candidate cases. Time window: 30\,min.

    \item \textbf{Evidence Agent}\,---\,Retrieves the top-50 supporting events per case, scored by a 4-factor relevance model:
    \begin{center}\small
    $\text{relevance} = 0.35 \times \text{temporal\_proximity} + 0.30 \times \text{IP\_match} + 0.15 \times \text{sensor\_diversity} + 0.20 \times \text{event\_type\_relevance}$
    \end{center}

    \item \textbf{Critic Agent}\,---\,Validates each case using a \textbf{5-factor confidence model}:
    \begin{center}\small
    $\text{confidence} = 0.25 \cdot S_{\text{det}} + 0.25 \cdot S_{\text{evid}} + 0.20 \cdot S_{\text{sensor}} + 0.15 \cdot S_{\text{temporal}} + 0.15 \cdot S_{\text{corr}}$
    \end{center}
    If validation fails (confidence\,$<$\,0.6 or insufficient evidence), the pipeline retries with expanded evidence search (up to 3 rounds).

    \item \textbf{Report Agent}\,---\,Generates a Markdown case report with executive summary, evidence table, timeline, detector reasoning, and defensive recommendations.
\end{enumerate}

\noindent\textbf{Terminology note:} ``Agent'' here refers to deterministic, rule-based pipeline stages---not autonomous AI agents with learned policies.

% ══════════════════════════════════════════════════════════════════════════
\section{Results}

\subsection{Benchmark Results (13 / 14 PCAPs Processed)}

\begin{table}[H]
\centering\small
\caption{Benchmark results across all processed PCAPs}
\label{tab:benchmark}
\begin{tabular}{@{} l c R{1cm} R{1.1cm} R{0.8cm} R{1.4cm} R{1.4cm} @{}}
\toprule
\textbf{PCAP} & \textbf{Label} & \textbf{Events} & \textbf{Detect.} & \textbf{Cases} & \textbf{Compr.} & \textbf{FP Proxy} \\
\midrule
synthetic-scan           & mal. & 868  & 1  & 1 & 1.0:1 & 52.3/hr \\
synthetic-beacon         & mal. & 352  & 1  & 1 & 1.0:1 & 2.8/hr  \\
synthetic-benign         & ben. & 168  & 0  & 0 & ---   & 0.0/hr  \\
synthetic-mixed          & mal. & 788  & 2  & 1 & 2.0:1 & 4.9/hr  \\
synthetic-dns-exfil      & mal. & 466  & 1  & 1 & 1.0:1 & 2.4/hr  \\
synthetic-multi-attacker & mal. & 805  & 2  & 2 & 1.0:1 & 8.5/hr  \\
synthetic-killchain      & mal. & 394  & 2  & 2 & 1.0:1 & 4.0/hr  \\
synthetic-slow-scan      & mal. & 2678 & 1  & 1 & 1.0:1 & 2.0/hr  \\
synthetic-noisy-benign   & ben. & 820  & 0  & 0 & ---   & 0.0/hr  \\
evasion-ip-rotation      & mal. & 552  & 0  & 0 & ---   & 0.0/hr  \\
evasion-jittery-beacon   & mal. & 390  & 1  & 1 & 1.0:1 & 3.0/hr  \\
evasion-slow-drip        & mal. & 1486 & 20 & 5 & 4.0:1 & 20.0/hr \\
wireshark-http           & ben. & 2    & 0  & 0 & ---   & 0.0/hr  \\
\bottomrule
\end{tabular}

\vspace{4pt}
\raggedright\footnotesize
CTU-13 Neris (56\,MB) timed out during processing and is excluded.  FP Proxy = detections per hour. Compression = detections\,:\,cases.
\end{table}

\subsection{Aggregate Metrics}

\begin{table}[H]
\centering\small
\caption{Key aggregate statistics}
\label{tab:aggregate}
\begin{tabular}{@{} l c @{}}
\toprule
\textbf{Metric} & \textbf{Value} \\
\midrule
Evidence completeness (detected malicious cases) & 100\% \\
False positives on benign traffic (3 PCAPs)      & 0 \\
Average compression ratio (malicious)             & 1.30\,:\,1 \\
Mean confidence (all runs)                        & 0.51 $\pm$ 0.36 \\
Effect size (Hedge's $g$)                         & \textbf{0.87} (large) \\
Sample sizes                                      & $n_{\text{mal}} = 10$,\; $n_{\text{ben}} = 3$ \\
\bottomrule
\end{tabular}

\vspace{4pt}
\raggedright\footnotesize
Hedge's $g$ is a small-sample-corrected effect size: $g = d \times \bigl(1 - \tfrac{3}{4(n_1+n_2)-9}\bigr)$. With $n_{\text{ben}}=3$ the confidence interval is wide; results are promising but statistically underpowered.
\end{table}

\subsection{Ablation Study (5 Configurations)}

All configurations tested on the \texttt{synthetic-multi-attacker} PCAP (805 events):

\begin{table}[H]
\centering\small
\caption{Ablation study results}
\label{tab:ablation}
\begin{tabular}{@{} L{3.6cm} C{1.3cm} C{1cm} C{1.3cm} C{1.4cm} C{1.4cm} @{}}
\toprule
\textbf{Configuration} & \textbf{Detect.} & \textbf{Cases} & \textbf{Compr.} & \textbf{Evidence} & \textbf{Confid.} \\
\midrule
(A) Suricata only        & 0 & 0 & 0:1   & 0\%   & 0.00 \\
(B) Detectors only       & 3 & 0 & 0:1   & 0\%   & 0.00 \\
(C) Full pipeline        & 3 & 2 & 1.5:1 & 100\% & \textbf{0.77} \\
(D) No critic agent      & 3 & 2 & 1.5:1 & 100\% & 0.92 \\
(E) Zeek sensor only     & 2 & 2 & 1.0:1 & 100\% & 0.67 \\
\bottomrule
\end{tabular}

\vspace{4pt}
\raggedright\footnotesize
Config A uses default Suricata community rules with no custom signatures. Its 0\% detection rate reflects the gap between signature-based and behavioral detection, not a flaw in Suricata itself.
\end{table}

\subsection{Adversarial Evasion Results}

\begin{table}[H]
\centering\small
\caption{Adversarial evasion test outcomes}
\label{tab:evasion}
\begin{tabularx}{\textwidth}{@{} l c c L{5.8cm} @{}}
\toprule
\textbf{Evasion Technique} & \textbf{Detected?} & \textbf{Cases} & \textbf{Analysis} \\
\midrule
IP rotation (5 sources)   & \textcolor{red}{\textbf{No}}  & 0 & Each IP stays below fan-out threshold (12\,<\,15). Reveals fundamental per-IP limitation. \\
Jittery DNS beacon        & \textcolor{green!50!black}{\textbf{Yes}} & 1 & Repeated-query signal strong enough despite high timing jitter. \\
Slow-drip exfiltration    & \textcolor{green!50!black}{\textbf{Yes}} & 5 & Domain-diversity and repeated-query signals compensate for lack of periodicity. 4:1 compression. \\
\bottomrule
\end{tabularx}
\end{table}

% ══════════════════════════════════════════════════════════════════════════
\section{Key Observations}

\begin{enumerate}[itemsep=4pt]
    \item \textbf{Multi-signal detection outperforms signatures.}\;
    Suricata alone (Config~A) detected 0 threats; the behavioral detectors found all but one evasion technique.

    \item \textbf{The critic agent improves calibration, not accuracy.}\;
    Config~D (no critic) scored 0.92 confidence---overconfident because it uses only 2 factors. The full 5-factor critic (Config~C) produced a better-calibrated 0.77 by penalizing low sensor diversity and absent cross-case correlation.

    \item \textbf{Multi-sensor fusion adds value.}\;
    Zeek-only (Config~E) missed 1 of 3 detections and scored 0.67 confidence vs.\ 0.77 with both sensors.

    \item \textbf{Agent orchestration creates actionable output.}\;
    Config~B produced 3 raw detections but 0 structured cases. The triage and evidence agents transform detections into analyst-ready packages.

    \item \textbf{Zero false positives on benign traffic.}\;
    Across 3 benign PCAPs (including a heavy-traffic stress test), the pipeline produced 0 spurious detections.

    \item \textbf{Evasion testing reveals honest limitations.}\;
    The IP rotation attack (5 sources, each scanning 12 targets) completely evaded detection. This exposes a fundamental architectural limitation: per-source-IP analysis cannot detect coordinated distributed campaigns.

    \item \textbf{Compression efficiency.}\;
    The slow-drip evasion PCAP generated 20 raw detections compressed into 5 cases (4:1 ratio), demonstrating meaningful noise reduction even on adversarial input.
\end{enumerate}

% ══════════════════════════════════════════════════════════════════════════
\section{Tools \& Technologies}

\begin{table}[H]
\centering\small
\begin{tabular}{@{} l l @{}}
\toprule
\textbf{Category} & \textbf{Tools} \\
\midrule
Network sensors   & Zeek, Suricata (Docker containers) \\
Language          & Python 3.13 \\
Core libraries    & pandas, NumPy, matplotlib, Scapy, PyYAML, Pydantic \\
Data format       & Apache Parquet \\
Infrastructure    & Docker Compose, Make (14 automation targets) \\
Testing           & pytest (51 unit tests), pytest-cov \\
\bottomrule
\end{tabular}
\end{table}

% ══════════════════════════════════════════════════════════════════════════
\section{Limitations \& Future Work}

\begin{itemize}[itemsep=3pt]
    \item \textbf{Synthetic-only evaluation.}\; All primary PCAPs are synthetically generated. Real production traffic is far noisier and more varied.
    \item \textbf{Small benign sample.}\; With only $n_{\text{ben}}=3$, the Hedge's~$g = 0.87$ has wide confidence intervals.
    \item \textbf{Two detection types only.}\; Lateral movement, non-DNS exfiltration, and APT-style attacks are not addressed.
    \item \textbf{No payload inspection.}\; The system operates entirely on flow metadata and DNS headers.
    \item \textbf{Batch mode only.}\; No real-time streaming capability; the 56\,MB CTU-13 PCAP exceeded the pipeline timeout.
    \item \textbf{Per-IP architecture gap.}\; Distributed attacks that keep each source below thresholds evade detection (demonstrated by the IP rotation evasion test).
\end{itemize}

\noindent\textbf{Future work:} Cross-source correlation (to defeat IP rotation), streaming ingestion, additional detection modules, and a pilot deployment on production SOC traffic to validate real-world performance.

% ══════════════════════════════════════════════════════════════════════════
\section{Conclusion}

This project demonstrates that a \textbf{deterministic, rule-based multi-agent pipeline}
can transform raw network telemetry into structured, evidence-grounded investigation
cases with transparent confidence scores. Across 13 processed PCAPs the system achieved
\textbf{100\% evidence completeness}, \textbf{zero false positives} on benign traffic,
and a \textbf{large effect size} (Hedge's~$g = 0.87$) separating malicious from benign
detection rates. The 5-configuration ablation study confirmed that each pipeline
component adds measurable value, with the critic agent providing better-calibrated
confidence (0.77 vs.\ 0.92 without it). Adversarial evasion testing honestly exposed
one architectural limitation (distributed IP rotation), while confirming detection
resilience against jittery beaconing and slow-drip exfiltration. The system is
reproducible (\texttt{make setup \&\& make benchmark}), well-tested (51 unit tests),
and designed with a clear path toward production SOC deployment.

\end{document}
